\chapter{Residual Networks and Skip Connections}

\section{The Degradation Problem}

As neural networks become deeper, one might expect their performance to improve due to increased expressivity. However, in practice, deeper networks can be harder to train.

\medskip

\noindent
\textbf{Observation.}  
Increasing depth may lead to higher training error, even though the deeper model can represent the same functions as a shallower one.

\medskip

\noindent
\textbf{Degradation problem.}
\begin{itemize}
    \item Deeper networks perform worse than shallower counterparts
    \item Not due to overfitting, but due to optimization difficulty
\end{itemize}

\medskip

\noindent
\textbf{Intuition.}  
In principle, additional layers could learn the identity mapping and preserve performance. However, standard training procedures struggle to discover such solutions.

\medskip

\noindent
\textbf{Connection to gradients.}
\begin{itemize}
    \item Deep compositions amplify vanishing and exploding gradients
    \item Information may degrade across layers
\end{itemize}

\medskip

\noindent
\textbf{Insight.}  
Depth increases expressivity, but also introduces optimization challenges that must be addressed through architectural design.

\section{Residual Learning}

Residual networks (ResNets) address the degradation problem by introducing \emph{skip connections}.

\medskip

\noindent
\textbf{Residual block.}
\[
z^{(\ell+1)} = z^{(\ell)} + F(z^{(\ell)}),
\]
where $F$ is a learnable transformation.

\medskip

\noindent
\textbf{Interpretation.}
\begin{itemize}
    \item Instead of learning a mapping $H(z)$, the network learns a residual:
    \[
    F(z) = H(z) - z
    \]
    \item The full mapping becomes $H(z) = z + F(z)$
\end{itemize}

\medskip

\noindent
\textbf{Advantages.}
\begin{itemize}
    \item Easier to learn small perturbations of identity
    \item Improves gradient flow across layers
    \item Reduces degradation
\end{itemize}

\medskip

\noindent
\textbf{Gradient propagation.}  
Gradients can flow directly through the identity connection:
\[
\frac{\partial z^{(\ell+1)}}{\partial z^{(\ell)}} = I + \frac{\partial F}{\partial z^{(\ell)}}.
\]

\medskip

\noindent
\textbf{Insight.}  
Skip connections provide alternative paths for information and gradients, stabilizing training.

\section{Dynamical Systems Perspective}

Residual networks admit an interpretation as discrete dynamical systems.

\medskip

\noindent
\textbf{Discrete dynamics.}
\[
z^{(\ell+1)} = z^{(\ell)} + F(z^{(\ell)}).
\]

\medskip

\noindent
\textbf{Continuous-time analogy.}  
This resembles a discretization of an ordinary differential equation:
\[
\frac{dz(t)}{dt} = F(z(t)).
\]

\medskip

\noindent
\textbf{Euler discretization.}  
With step size $\Delta t = 1$,
\[
z(t + \Delta t) \approx z(t) + \Delta t \, F(z(t)).
\]

\medskip

\noindent
\textbf{Interpretation.}
\begin{itemize}
    \item Layers correspond to time steps
    \item Network depth corresponds to evolution over time
\end{itemize}

\medskip

\noindent
\textbf{Stability.}  
Residual connections help maintain stable trajectories in representation space.

\medskip

\noindent
\textbf{Insight.}  
Viewing neural networks as dynamical systems provides a framework for understanding depth, stability, and training dynamics.

\section{Very Deep Architectures}

Residual connections enable the training of extremely deep networks.

\medskip

\noindent
\textbf{Empirical observation.}
\begin{itemize}
    \item Networks with hundreds or thousands of layers can be trained successfully
    \item Performance improves with depth when using residual connections
\end{itemize}

\medskip

\noindent
\textbf{Architectural patterns.}
\begin{itemize}
    \item Stacking residual blocks
    \item Combining with normalization layers
    \item Using bottleneck structures for efficiency
\end{itemize}

\medskip

\noindent
\textbf{Scaling behavior.}
\begin{itemize}
    \item Deeper models capture more complex patterns
    \item Training remains stable due to skip connections
\end{itemize}

\medskip

\noindent
\textbf{Generalization.}  
Residual networks often generalize well despite large capacity.

\medskip

\noindent
\textbf{Limitations.}
\begin{itemize}
    \item Increased computational cost
    \item Architectural complexity
\end{itemize}

\medskip

\noindent
\textbf{Insight.}  
Careful architectural design makes depth practically usable.

\section{A Unifying Perspective}

Residual networks illustrate how architectural innovations can address fundamental challenges in deep learning.

\medskip

\noindent
\textbf{Key components.}
\begin{itemize}
    \item Skip connections
    \item Residual mappings
    \item Stable gradient flow
\end{itemize}

\medskip

\noindent
\textbf{Core ideas.}
\begin{itemize}
    \item Learn perturbations rather than full transformations
    \item Facilitate optimization through identity mappings
    \item Interpret networks as dynamical systems
\end{itemize}

\medskip

\noindent
\textbf{Key insight.}  
Architectural design can fundamentally change the behavior of learning algorithms.

\section{Outlook}

This chapter introduced residual networks and skip connections:
\begin{itemize}
    \item the degradation problem in deep networks,
    \item residual learning as a solution,
    \item connections to dynamical systems,
    \item enabling very deep architectures.
\end{itemize}

\medskip

In the next chapter, we study convolutional neural networks, which exploit spatial structure in data.

\medskip

\noindent
\textbf{Guiding principle.}  
Effective architectures combine mathematical structure with practical considerations to make deep learning feasible at scale.